% UNREVIEWED DRAFT: written during checkpoint05 freeze.
% Not part of the results inventory until independent root audit.
\paragraph*{Unreviewed draft, not in the results inventory.}

\begin{saresult}{B.047}{Proved}{Completion with an arbitrary rational coefficient overring}
In the setting of B.044, let $R\subset L\subset K$ be
any overring and put $B=L\otimes_R A\subset Q$.
For every $n$ there is a natural isomorphism
\[
 \Omega_n(R,L)\otimes_R A\simeq\Omega_n(A,B).
\]
Here $\Omega_n(A,B)$ denotes the polynomials in $B[X_1,
\ldots,X_n]$ taking $A^n$ into $A$. The tensor maps
correspond, and existence of the coaddition and
comultiplication expressions passes to this base change.
\end{saresult}
\begin{proof}
Flatness embeds $B=L\otimes_R A$ in $K\otimes_R A=Q$.
Within a fixed finite rectangular degree bound with $N$
monomials, the coefficient module of $\Omega_n(R,L)$ is
the intersection of the bounded-degree integer-valued
module with $L^N$, inside $K^N$. Flat base change preserves
this intersection, by tensoring the kernel sequence for
the difference map from the direct sum of the two
submodules to $K^N$. B.044 identifies the first module
after base change, and the second becomes $B^N$.
Take the filtered union over degree bounds. These
identifications preserve multiplication in disjoint
variables and polynomial substitution. Finally every
element of the algebraic tensor product is a finite
$A$-linear combination of original polynomials. The
expressions for $F(X+Y)$ and $F(XY)$ therefore extend
by this linearity. No finite generation of $L$ is used.
\end{proof}

\begin{saresult}{B.048}{Proved}{Finite-rank modules over a complete discrete valuation ring}
Let $V$ be a complete discrete valuation ring with
uniformizer $t$ and fraction field $F$. Let $W$ be a
finite-dimensional $F$-vector space, and let $M\subset W$
be a $V$-submodule containing a full $V$-lattice.
There is an $F$-vector subspace
$D=\bigcap_{j\ge0}t^jM$ such that $M/D$ is a finite
free $V$-module. More precisely, after a choice of
$F$-basis, $M$ is a direct sum of finitely many
fractional ideals $t^aV$ and finitely many copies of $F$.
\end{saresult}
\begin{proof}
Choose a full lattice $\Lambda\subset M$ and identify
$(W,\Lambda)$ with $(F^d,V^d)$. Put $E=F/V$.
We first record the elementary duality facts used below.
The divisible module $E$ is injective over the principal
ideal domain $V$. For a direct proof of the extension
property, extend a homomorphism maximally by Zorn's
lemma. To adjoin a new element $x$, its annihilator
relative to the old submodule is an ideal $(a)$.
If $a\ne0$, divisibility solves $ay=f(ax)$; if the
ideal is zero choose any $y$. Sending $x$ to $y$
extends the map, a contradiction unless the domain was
already the whole module.
Also $E$ detects every nonzero torsion module: a nonzero
element has annihilator $(t^s)$, its cyclic submodule
maps nontrivially to $t^{-s}V/V$, and injectivity extends
this map to the whole module.

An endomorphism of $E$ is multiplication by a unique
element of $V$. Indeed its value on $t^{-s}+V$ is
$a_s t^{-s}+V$ for $a_s\in V/t^sV$, and compatibility
is exactly $a_{s+1}\equiv a_s\pmod {t^s}$. Completeness
identifies the compatible sequence with an element of
$V$. Consequently every homomorphism $E^d\to E$
is induced by a vector in $V^d$.

Let
\[
 M^\#=\{\phi\in W^*: \phi(M)\subset V\}
                                      \subset\Lambda^*.
\]
Then
\[
 M=\{x\in W: \phi(x)\in V\text{ for every }
                                      \phi\in M^\#\}.       \tag{B48.1}
\]
Only the reverse inclusion needs proof. If $x\notin M$,
the nonzero torsion element $x+M$ in $W/M$ is detected
by a homomorphism to $E$. Composing with
$W/\Lambda\simeq E^d$ and using the preceding
endomorphism computation gives an $F$-linear form
$\phi$ integral on $M$ but not on $x$.

The submodule $M^\#\subset\Lambda^*$ is finite free,
since $V$ is Noetherian. Smith reduction of its matrix
inside the finite free module $\Lambda^*$ puts it in
the form
$t^{a_1}Ve_1^*\oplus\cdots\oplus t^{a_s}Ve_s^*$
with zero remaining coordinates. Formula (B48.1)
then gives
\[
 M=t^{-a_1}Ve_1\oplus\cdots\oplus t^{-a_s}Ve_s
           \oplus Fe_{s+1}\oplus\cdots\oplus Fe_d.
\]
Its intersection of all $t^j$ multiples is exactly the
last vector-space summand, proving every assertion.
\end{proof}

\begin{saresult}{B.049}{Proved}{Epimorphic overrings of a complete one-dimensional order}
Let $V$ be a complete discrete valuation ring with finite
residue field, let $A$ be a local commutative $V$-algebra
finite and torsion-free over $V$, and put
$Q=A\otimes_V\operatorname{Frac}(V)$.
For any overring $A\subset B\subset Q$, the multiplication
map $B\otimes_A B\to B$ is an isomorphism if and only if
\[
                            B=A+eQ
\]
for an idempotent $e\in Q$. Moreover injectivity of the
second polynomial tensor map for $\Omega(A,B)$ forces
this ring-epimorphism condition.
\end{saresult}
\begin{proof}
For the last assertion choose any $b\in B$. There is
a nonzerodivisor $d\in A$ with $db\in A$; a power
of the uniformizer of $V$ suffices. The finite ring
$A/dA$ has a monic null polynomial, for instance the
product of $X-a$ over a complete set of representatives.
Call it $P$. Then $P,bP\in\Omega(A,B)$. The tensor
\[
                       bP\otimes P-P\otimes bP
\]
has zero polynomial image, hence is zero if the tensor
map is injective. Taking the coefficient of degree
$\deg P$ in both factors gives
$b\otimes1=1\otimes b$ in $B\otimes_A B$.
This for every $b$ makes multiplication an isomorphism,
with inverse $b\mapsto b\otimes1$.

Now assume multiplication is an isomorphism.
Apply B.048 to $B\subset Q$, using the full finite
free lattice $A$. The divisible part
$D=\bigcap_j t^jB$ is an ideal of $B$ and a vector
space over $\operatorname{Frac}(V)$. It is in fact
an ideal of $Q$, since $Q=A[1/t]$. The quotient
$C=B/D$ is finite over $V$, hence finite over $A$.
The composite $A\to C$ is a ring epimorphism because
$A\to B$ and $B\to C$ are. If $C\ne0$, its fiber
$C/\mathfrak m_A C$ is nonzero by Nakayama and is a
finite-dimensional algebra over $k=A/\mathfrak m_A$.
The multiplication isomorphism on this fiber implies
$(\dim_k C/\mathfrak m_A C)^2=\dim_k C/\mathfrak m_A C$,
so the dimension is one and the unit generates it.
Nakayama applied to $C/\operatorname{im}(A)$ shows
that $A\to C$ is surjective. The same conclusion is
automatic when $C=0$. Therefore $B=A+D$.

Put $I=A\cap D$. It is a finite $V$-module. From
$B/A\simeq D/I$ and the epimorphism condition,
\[
              B\otimes_A(D/I)=B\otimes_A(B/A)=0.
\]
As $D$ is already a $Q$-module,
$B\otimes_A D\simeq Q\otimes_QD=D$.
Right exactness now gives $D=BI=I+DI$.
The quotient $D/DI$ is a vector space over
$\operatorname{Frac}(V)$ generated as a $V$-module
by the image of the finite module $I$. Such a vector
space must be zero: multiplication by $t$ is
surjective, and Nakayama applies to a finite
$V$-module. Hence $D=DI\subset D^2\subset D$.
An idempotent ideal in the Artinian ring $Q$ is
$eQ$: in each local Artinian factor a proper ideal
is nilpotent, so an idempotent ideal is either zero
or the whole factor. This proves necessity.

Conversely, for $B=A+eQ$, write $f=1-e$. Then
$B=eQ\times fA$. The map $A\to fA$ is a quotient,
so $fA\otimes_A fA\simeq fA$. Localization gives
$eQ\otimes_A eQ\simeq eQ$, and the cross terms
$eQ\otimes_A fA$ and $fA\otimes_A eQ$ vanish,
because after inverting $t$ they involve the
orthogonal factors $eQ$ and $fQ$. Multiplication
is therefore an isomorphism on the whole product.
\end{proof}

\begin{saresult}{B.050}{Proved}{All one-dimensional Noetherian domains and every torsion-free coefficient extension}
Let $R$ be an infinite one-dimensional Noetherian domain,
let $K=\operatorname{Frac}(R)$, and let $R\subset S$
be an $R$-torsion-free unital extension. Put
$L=S_{\mathrm{rat}}\subset K$ as in B.014.
The coaddition and comultiplication tensors are
uniquely determined if and only if at each maximal
ideal $\mathfrak m$ with finite residue field either
\[
 L_{\mathfrak m}=R_{\mathfrak m},
 \qquad\hbox{or}\qquad
 L_{\mathfrak m}=K\ \hbox{and}\
 \widehat{R_{\mathfrak m}}\hbox{ is a domain}.              \tag{B50.1}
\]
Equivalently, in the second alternative the local
normalization is finite and local. In every positive
case all tensor maps are isomorphisms and the
compatible plethory is unique. There is no
finite-normalization assumption on $R$.
\end{saresult}
\begin{proof}
B.014 reduces to $L$ and B.033 localizes the question.
At infinite-residue maximal ideals Vandermonde
interpolation gives the ordinary polynomial ring.
Fix a finite-residue local ring $R$, put
$A=\widehat R$, $Q=K\otimes_R A$, and
$B=L\otimes_R A\subset Q$. B.047 transfers the
question to $\Omega(A,B)$. The completion $A$ is
one-dimensional Cohen--Macaulay. As in the first
paragraphs of B.045, it is finite and torsion-free
over a complete discrete valuation subring with
finite residue field.

If the tensor map is injective, B.049 forces
$B=A+eQ$ for an idempotent $e$. If $e=0$, then
$B=A$ and faithful flatness gives $L=R$.
Suppose $e\ne0$. If any selected local Artinian
factor $Q_0$ is nonreduced, the two polynomials
$F_r,H_r$ constructed in B.045 have all coefficients
in that selected factor $e_0Q\subset eQ\subset B$.
They therefore belong to $\Omega(A,B)$, and the same
balanced functional proves a nonzero tensor kernel
there. Thus every selected factor must be a field.

Suppose there is some other generic factor besides
a selected field $Q_0$, and let $e_0$ be its
idempotent. Let $V\subset Q_0$ be the integral
closure of the complete discrete valuation subring.
It is a complete discrete valuation ring with finite
residue field, and the projection of $A$ lies in $V$.
Choose its uniformizer $u$ and $c\ge1$ with
$u^c e_0V\subset A$. Put $x=u^c e_0\in A$
and $r=c+1$.
The ideal $I_0=A\cap e_0Q$ is proper, since
$e_0\ne1$, and hence lies in $\mathfrak m_A$.
For every $h\in\Omega(A,B)$ the difference
$h(x)-h(0)$ has zero components outside $Q_0$;
as both values belong to $A$, it lies in $I_0$.

There is an integer-valued polynomial $H$ on $V$
such that $H(u^{cr})-H(0)$ is a unit. For example
the finite congruence-class indicator construction
in B.024 separates $0$ and $u^{cr}$ modulo a power
$u^M$ with $M>cr$. Define the polynomial, supported
on $Q_0$,
\[
                    G=u^c e_0\bigl(H-H(0)\bigr).
\]
Its coefficients lie in $e_0Q\subset B$, it takes
$A$ into $u^ce_0V\subset A$, and $G(0)=0$.
Moreover its selected valuation at $x^r$ is exactly
$c$. If multiplicative tensors existed for every
polynomial in $\Omega(A,B)$, iterating the binary
expressions would give an $r$-fold expression for
$G(X_1\cdots X_r)$. Apply the difference of evaluation
at $x$ and at zero in each variable. Every term on
the tensor side is in $\mathfrak m_A^r$, whereas
the other side is $G(x^r)$.
The projection of $\mathfrak m_A$ lies in the
maximal ideal of $V$, since this valuation ring
dominates the selected reduced local branch.
Thus $\mathfrak m_A^r$ has valuation at least
$r=c+1$, contradicting the value $c$ just found.
So this case fails even existence.

The only remaining possibility with $e\ne0$ is
that $Q$ itself is one field and $B=Q$.
Because $A\hookrightarrow Q$, this is exactly
the case that $A$ is a domain and $B=Q$.
Faithful flatness applied to $K/L$ gives $L=K$.
This proves necessity of (B50.1).

Conversely $L=R$ gives the ordinary polynomial ring.
If $L=K$ and the completion is a domain, B.046
gives all tensor isomorphisms. The same theorem
identifies this hypothesis with finite local
normalization. Applying these local conclusions
at every maximal ideal proves sufficiency globally
by testing kernels and cokernels. The unique
plethory compatibility is B.009.
\end{proof}
